% Strategic Management – Week 9
% Compile: pdflatex starbucks-case.tex

\documentclass[10pt,a4paper]{article}
\usepackage[utf8]{inputenc}
\usepackage[T1]{fontenc}
\usepackage[margin=1.2cm]{geometry}
\usepackage{enumitem}
\setlist{nosep,leftmargin=*}
\usepackage{xcolor}
\usepackage{marginnote}
\usepackage{fancyhdr}
\usepackage{microtype}

\definecolor{eublue}{RGB}{0,51,153}
\definecolor{notegrey}{RGB}{90,90,90}
\pagestyle{fancy}
\fancyhf{}
\fancyfoot[L]{\small Strategic Mgmt | Week 9}
\fancyfoot[R]{\small\thepage}
\raggedright

\newcommand{\handnote}[1]{\marginnote{\raggedright\footnotesize\color{notegrey}\textit{#1}\vspace{0.4em}}}
\newcommand{\hlpink}[1]{\colorbox{pink!25}{\textbf{#1}}}

\begin{document}
\begin{center}
{\large\bfseries\color{eublue} Strategic Management -- Starbucks Case}\\[0.2em]
\footnotesize Mar 23, 2026 (Week 9). Break-down of what is going on + how to argue it in class.
\end{center}
\vspace{0.6em}

\section*{\color{eublue}1. Case in one paragraph (what is going on)}
Starbucks is used as a ``live'' case for how a firm turns an initial popular offer into a repeatable system that can be scaled across locations and time. The key strategic tension is not only ``how did growth happen?'' but also ``what lets the firm keep its advantage after competitors react and markets change?'' The case points to sustained success coming from a business model that stays consistent in the customer value delivered, while still adapting continuously in execution (operations, experience, and fit to markets).

\section*{\color{eublue}2. Strategy breakdown (mechanisms)}
\begin{enumerate}
  \item \textbf{Repeatable business model execution.}
  \begin{itemize}
    \item Store-level operations are designed so that the customer experience is recognizable across locations.
    \item Brand consistency and product/service delivery act like a ``system'' rather than a collection of isolated moves.
  \end{itemize}

  \item \textbf{Growth logic: scale without losing identity.}
  \begin{itemize}
    \item Expansion works when the company can replicate not just the stores, but the routines behind them (training, standards, and coordination).
    \item The advantage is reinforced as learning accumulates through operations at scale.
  \end{itemize}

  \item \textbf{Continuous adaptation (not one-time success).}
  \begin{itemize}
    \item The firm sustains growth by updating what it offers and how it delivers value as customer behavior and competitive conditions evolve.
    \item Adaptation is framed as changing ``implementation'' while protecting the core differentiators (what customers come for).
  \end{itemize}

  \item \textbf{Reinvention / turnaround logic (when performance is threatened).}
  \begin{itemize}
    \item When the original momentum is challenged, the strategy focuses on rebuilding store execution and restoring the customer experience.
    \item The lesson: advantage renewal requires active management choices, not automatic defense.
  \end{itemize}

  \item \textbf{International scaling logic.}
  \begin{itemize}
    \item The brand and system travel, but local markets require adaptation in how the offer is positioned and executed.
    \item Sustained growth depends on turning local expansion into repeatable learning that strengthens execution over time.
  \end{itemize}
 \end{enumerate}

\section*{\color{eublue}3. Map the case to course frameworks (how to argue)}
\subsection*{\color{eublue}A. RBV / VRIO (internal explanation)}
\begin{enumerate}
  \item \textbf{Value (V).} The Starbucks system delivers recognizable customer value (experience + consistency), which supports demand and long-term profitability.
  \item \textbf{Rarity (R).} It is not only ``a coffee shop''; it is the combination of brand, operations, and customer-experience routines that is harder to find as a single package.
  \item \textbf{Inimitability (I).} Competitors can imitate products quickly, but copying the full operating system (training + routines + standards + brand experience) is slower and culturally embedded.
  \item \textbf{Non-substitutability (N).} Customers may substitute coffee, but it is harder to substitute the specific ``repeatable experience'' with a simple alternative without losing what they value.
  \item \textbf{Organization (O).} The firm must be organized to exploit the resources (execution standards, capability building, and management discipline across locations).
\end{enumerate}

\subsection*{\color{eublue}B. Competitive advantage as a renewal problem}
The case pushes the idea that competitive advantage must be renewed: rivalry and imitation gradually erode early wins, so sustained advantage depends on continual improvements while protecting core differentiators.

\subsection*{\color{eublue}C. Value chain (where the execution happens)}
\begin{itemize}
  \item \textbf{Inbound logistics:} ensure consistent input quality (coffee sourcing and standards).
  \item \textbf{Operations:} turn the brand promise into store routines (how the product and service are delivered).
  \item \textbf{Marketing \& sales / service:} reinforce the experience and make the offering recognizable.
  \item \textbf{Supporting activities:} HR/training and coordination capabilities that scale the system across geographies.
\end{itemize}

\section*{\color{eublue}4. Risks and counter-arguments (what could break the advantage)}
\begin{enumerate}
  \item \textbf{Commoditization of products.} If coffee itself becomes easier to match, differentiation must shift to experience and execution.
  \item \textbf{Experience drift at scale.} Expansion can create inconsistency if training/standards weaken.
  \item \textbf{Competitive response.} Rivals can copy surface features; sustained advantage requires deeper operational coherence.
  \item \textbf{Market saturation / cost pressure.} Growth may slow if expansion opportunities diminish or costs rise.
\end{enumerate}

\section*{\color{eublue}5. What to say in class (speaking templates)}
\subsection*{\color{eublue}A. 60--90 seconds (oral answer skeleton)}
\begin{enumerate}
  \item \textbf{Opening.} Starbucks shows how a firm turns initial popularity into sustainable advantage through positioning, execution, and scaling.
  \item \textbf{Why it worked.} The business model was repeatable across locations, so customers received a consistent experience; operations and experience were managed as a system.
  \item \textbf{Why growth was sustained.} Growth came from continuous adaptation while protecting core differentiators; advantage had to be renewed, not only defended.
  \item \textbf{Close.} The main lesson is that competitive advantage must be actively renewed over time.
\end{enumerate}

\subsection*{\color{eublue}B. 30 seconds cold-call version (3 lines)}
\begin{enumerate}
  \item \textbf{Topic:} Starbucks shows how initial popularity becomes sustainable advantage.
  \item \textbf{Logic:} Repeatable business model + system execution + continuous adaptation.
  \item \textbf{Action/implication:} Advantage persists only if the firm renews execution as markets change.
\end{enumerate}

\section*{\color{eublue}6. Answers to the case questions (ready-to-say)}
% Note: wording of your case synopsis may vary slightly by edition. This matches the typical 3-part prompt shown in your `starbucks-1.JPG`.
% Key questions (as visible in `starbucks-1.JPG`):
% - "Why has Starbucks' strategy succeeded so well?"
% - "Is Starbucks' competitive advantage sustainable?"
% - "Give your reason."
% - Advice prompt: what Mr. Schultz should amend/refine in order to remain successful.

\subsection*{\color{eublue}Key quotes / phrases from your pages}
\begin{enumerate}
  \item ``The future''
  \item ``Rekindling growth''
  \item ``Crisis and renaissance''
  \item ``Acquiring coffee beans of a high, consistent quality''
  \item ``Employee involvement''
  \item ``Closing Case''
\end{enumerate}

\subsection*{\color{eublue}Q1. Why has Starbucks' strategy succeeded so well?}
Answer (based on your notes + what the reading labels):
\\
``POPULARITY BUILT THROUGH REPEATABLE BUSINESS MODEL EXECUTION.''
\\
``OPERATIONS + CUSTOMER EXPERIENCE WERE MANAGED AS A SYSTEM, NOT AS ISOLATED MOVES.''
\\
``GROWTH SUSTAINED THROUGH CONTINUOUS ADAPTATION, NOT ONE-TIME MOVE.''
\\
\\
\textbf{Unit 5 (Competitive Strategy) 5.4 alignment (framework lines):}
\\
``PRODUCT DIFFERENTIATION''
\\
``Product diff ≠ Price Premium''
\\
``Costs NOT ALWAYS → C.A.''
\\
\textbf{Proof from class notes (`2026-03-23-class-notes.md`):}
\\
``POPULARITY BUILT THROUGH REPEATABLE BUSINESS MODEL EXECUTION'' (II.A.1)
\\
``OPERATIONS + CUSTOMER EXPERIENCE WERE MANAGED AS A SYSTEM, NOT AS ISOLATED MOVES'' (IV.B.2)
\\
``GROWTH SUSTAINED THROUGH CONTINUOUS ADAPTATION, NOT ONE-TIME MOVE'' (II.A.2)
\\
\\
\textbf{Proof from reading pages (`starbucks-*.JPG`):}
\\
``The future'' (starbucks-2)
\\
``Closing Case'' (starbucks-5)
\\
``Acquiring coffee beans of a high, consistent quality'' (starbucks-5)
\\
``Employee involvement'' (starbucks-5)

\textbf{Proof from `sm-unit-5-competitive-strategy.pdf` (Unit 5, 5.4):}
\\
``PRODUCT DIFFERENTIATION''
\\
``Product diff ≠ Price Premium''
\\
``Costs NOT ALWAYS → C.A.''
\\
``Stuck in the “middle”?'' 

\subsection*{\color{eublue}Q2. Is Starbucks' competitive advantage sustainable? Why or why not?}
Answer (based on your notes + what the reading labels):
\\
``THE MAIN LESSON IS THAT COMPETITIVE ADVANTAGE MUST BE RENEWED, NOT JUST DEFENDED.''
\\
\textbf{Unit 5 (Competitive Strategy) 5.4 support (framework lines):}
\\
``Costs NOT ALWAYS → C.A.''
\\
``Stuck in the “middle”?'' 
\\
\\
\textbf{Proof from class notes (`2026-03-23-class-notes.md`):}
\\
``THE MAIN LESSON IS THAT COMPETITIVE ADVANTAGE MUST BE RENEWED, NOT JUST DEFENDED'' (IV.D.1)
\\
\\
\textbf{Proof from reading pages (`starbucks-*.JPG`):}
\\
``Crisis and renaissance'' (starbucks-4)
\\
``Rekindling growth'' (starbucks-3)
\\
\\
\textbf{Proof from `sm-unit-5-competitive-strategy.pdf` (Unit 5, 5.4):}
\\
``Costs NOT ALWAYS → C.A.''
\\
``Stuck in the “middle”?'' 

\subsection*{\color{eublue}Q3. If Starbucks wants to stay successful, what should it focus on next?}
Answer (next focus), using your notes as the “what to do”:
\\
``REVIEW FRAMEWORKS FOR SUSTAINING COMPETITIVE ADVANTAGE.''
\\
``PREPARE TO MAP CASE OBSERVATIONS TO COURSE MODELS.''
\\
\\
\textbf{Unit 6 (Corporate Strategy) 6.2 support (framework lines):}
\\
``Try to consolidate its position in the markets and strengthen its competitive position in the industry (lower costs or more value added).''
\\
``Considered as low risk strategy “always doing the same products for the same customers needs”''
\\
\textbf{Proof from class notes (`2026-03-23-class-notes.md`):}
\\
``REVIEW FRAMEWORKS FOR SUSTAINING COMPETITIVE ADVANTAGE'' (III.B.1)
\\
``PREPARE TO MAP CASE OBSERVATIONS TO COURSE MODELS'' (III.B.2)
\\
\\
\textbf{Proof from reading pages (`starbucks-*.JPG`):}
\\
``Closing Case'' (starbucks-5)
\\
``Acquiring coffee beans of a high, consistent quality'' (starbucks-5)
\\
``Employee involvement'' (starbucks-5)
\\
``Rekindling growth'' (starbucks-3)
\\
\\
\textbf{Proof from `sm-unit-6-corporate-strategy-development.pdf` (Unit 6, 6.2):}
\\
``Try to consolidate its position in the markets and strengthen its competitive position in the industry (lower costs or more value added).''
\\
``Considered as low risk strategy “always doing the same products for the same customers needs”''

\section*{\color{eublue}Transcription from images (best-effort, legible parts only)}
\subsection*{\color{eublue}A. Case question box (from `starbucks-1.JPG`)}
\begin{quote}
\hlpink{Case question / Synopsis}

\handnote{Use these exact prompts for your Q1/Q2/Q3 answers.}

\begin{enumerate}
  \item \hlpink{Why has Starbucks' strategy succeeded so well?} \handnote{Class proof: ``POPULARITY BUILT THROUGH REPEATABLE BUSINESS MODEL EXECUTION.'' (II.A.1)}
  \item \hlpink{Is Starbucks' competitive advantage sustainable? Give your reason.} \handnote{Class proof: ``THE MAIN LESSON IS THAT COMPETITIVE ADVANTAGE MUST BE RENEWED, NOT JUST DEFENDED.'' (IV.D.1)}
  \item \hlpink{Advise Mr. Schultz on whether and how Starbucks should amend its strategy in order to remain successful.} \handnote{Class proof: ``REVIEW FRAMEWORKS FOR SUSTAINING COMPETITIVE ADVANTAGE.'' (III.B.1)}
\end{enumerate}
\end{quote}

\subsection*{\color{eublue}B. Closing Case bullets (from `starbucks-5.JPG`)}
\begin{quote}
\textbf{Closing Case / Starbucks Corporation} \handnote{Use this page to ground execution details.}

\begin{itemize}
  \item \hlpink{Acquiring coffee beans of a high, consistent quality} \handnote{Link to execution foundations: class notes focus on how advantage is maintained after success.}
  \item \hlpink{Employee involvement} \handnote{Link to sustaining the strategy over time via people/experience.}
\end{itemize}

% Body paragraphs: too small in these photos to transcribe verbatim without risking mistakes.
\textit{[Full body paragraphs from `starbucks-5.JPG` are marked unreadable/too blurry here; send zoomed crops if you want verbatim transcription.]}
\end{quote}

\end{document}

